% =====================================================================
%  4.4 Drøfting av funksjoner — ekstra utfordring
% =====================================================================

\begin{infoboks}{Velg metode selv}
Arket bygger på 4.4. Vis regning og begrunnelse. Bruk eksakte svar når oppgaven ikke ber om avrunding.
Ved optimering bruker du de fire stegene og svarer med en setning og enhet.
\end{infoboks}

\nivaahode{2}

\oppg{U1}{2}{\uh}
La $f(x) = x^3 - 3x + 2$.
\begin{tasks}(1)
  \task Vis at $x = 1$ er et nullpunkt. Bruk polynomdivisjon og skriv $f$ på faktorisert form.
  \task Drøft $f$: fortegnslinje for $f'$, ekstremalpunkter og skisse.
  \task Nullpunktet $x = 1$ er et dobbelt nullpunkt. Hva forteller det om grafen der? Sammenlign med ekstremalpunktene.
\end{tasks}

\oppg{U2}{2}{\mh}
En bedrift selger $x$ enheter. Fortjenesten er $F(x) = -0,02x^2 + 40x - 8000$ kroner, for $0 \le x \le 1500$ (tenkt modell).
\begin{tasks}(1)
  \task Bruk de fire stegene til å finne hvor mange enheter som gir størst fortjeneste. Svar i en setning med enhet.
  \task Vurder modellen: Er det realistisk at $F$ er en andregradsfunksjon? Hva kan gå galt for store $x$?
\end{tasks}

\oppg{U3}{2}{\uh}
Et rektangel har to hjørner på $x$-aksen og to hjørner på parabelen $y = 12 - x^2$, symmetrisk om $y$-aksen. Rektangelet har bredde $2x$.
Bruk de fire stegene til å finne størst mulig areal. Oppgi sidene i rektangelet.

\oppg{U4}{2}{\prog}
Programmet prøver $x$-verdier med steg $0,01$ og finner størst volum for esken fra eksempel 2.
\begin{kode}
def V(x):
    return x * (20 - 2 * x) * (30 - 2 * x)

beste_x = 0
beste_V = 0
for i in range(0, 1001):
    x = i / 100
    if V(x) > beste_V:
        beste_V = V(x)
        beste_x = x
print(round(beste_x, 2), round(beste_V, 4))
\end{kode}
\begin{tasks}(1)
  \task Kjør programmet. Hva skriver det ut? Hvorfor går \texttt{i} fra $0$ til $1000$?
  \task Sammenlign med det eksakte svaret fra $V'(x) = 0$. Hvor stor er feilen i $x$ og i $V$? Forklar hvorfor feilen i $V$ er mye mindre.
  \task Endre programmet til steg $0,001$. Da må du endre \texttt{range}, \texttt{i / 100} og \texttt{round(beste\_x, 2)}. Hva skriver det ut nå?
\end{tasks}

\nivaahode{3}

\oppg{U5}{3}{\uh}
La $f(x) = x^3 + ax$, der $a$ er et tall.
\begin{tasks}(1)
  \task For hvilke $a$ har $f$ ingen topp- eller bunnpunkt? Begrunn med $f'$.
  \task For hvilken $a$ har $f$ et terrassepunkt? Hvor ligger det?
\end{tasks}

\oppg{U6}{3}{\uh}
Bevis at et tredjegradspolynom $f(x) = px^3 + qx^2 + rx + s$ ($p \ne 0$) har enten $0$ eller $2$ ekstremalpunkter.
\hint{Se på $f'$ og antall nullpunkter. Hva avgjør om fortegnet skifter?}

\oppg{U7}{3}{\uh}
Finn tredjegradspolynomet med toppunkt $(-1, 7)$ og bunnpunkt $(1, 3)$. Finnes det flere slike polynomer? Begrunn.

\begin{oppsummering}
Fortegnslinjen for $f'$ forteller hvor $f$ stiger og synker. Et nullpunkt for $f'$ er bare en kandidat.
Største og minste verdi på et lukket intervall finnes ved å sammenligne endepunkter og kandidater.
Optimering gjøres i fire steg og avsluttes med en setning med enhet.
\end{oppsummering}
